\pagebreak

\section{Intercollegiate Math Tournament (ICMT) Integration Bee}
\subsection*{2026}

\subsubsection*{Regular Round: Bracket 1}
\begin{questions}

\question \[\int_0^4 |x^2-4|\,dx\]
\begin{solution}
Split the domain where $x^2-4$ changes sign, at $x=2$:
\[
\int_0^4 |x^2-4|\,dx=\int_0^2(4-x^2)\,dx+\int_2^4(x^2-4)\,dx.
\]
The first piece is $\left[4x-\dfrac{x^3}{3}\right]_0^2=8-\dfrac{8}{3}=\dfrac{16}{3}$.
The second piece is $\left[\dfrac{x^3}{3}-4x\right]_2^4=\left(\dfrac{64}{3}-16\right)-\left(\dfrac{8}{3}-8\right)=\dfrac{16}{3}+\dfrac{16}{3}=\dfrac{32}{3}$.
Adding, $\dfrac{16}{3}+\dfrac{32}{3}=16$.
\end{solution}

\question \[\int_0^1 x^3\ln(x)\,dx\]
\begin{solution}
Integrate by parts with $u=\ln x,\ dv=x^3\,dx$, so $du=\dfrac{dx}{x},\ v=\dfrac{x^4}{4}$:
\[
\int_0^1 x^3\ln x\,dx=\left[\frac{x^4}{4}\ln x\right]_0^1-\int_0^1\frac{x^3}{4}\,dx.
\]
The boundary term vanishes at both ends ($x^4\ln x\to0$ as $x\to0^+$). Thus
\[
=-\frac14\left[\frac{x^4}{4}\right]_0^1=-\frac{1}{16}.
\]
\end{solution}

\question \[\int_0^1 \frac{dx}{\sqrt{x}\sqrt{1-x}}\]
\begin{solution}
Substitute $x=\sin^2\theta$, $dx=2\sin\theta\cos\theta\,d\theta$, $\sqrt{x}=\sin\theta$, $\sqrt{1-x}=\cos\theta$:
\[
\int_0^{\pi/2}\frac{2\sin\theta\cos\theta}{\sin\theta\cos\theta}\,d\theta=\int_0^{\pi/2}2\,d\theta=\pi.
\]
\end{solution}

\end{questions}

\subsubsection*{Regular Round: Bracket 2}
\begin{questions}

\question \[\int_0^{\pi/3}\frac{\sec(x)\tan(x)}{\sec(x)(\cos(x)+1)}\,dx\]
\begin{solution}
Cancel $\sec x$ and write $\tan x=\dfrac{\sin x}{\cos x}$:
\[
\int_0^{\pi/3}\frac{\tan x}{\cos x+1}\,dx=\int_0^{\pi/3}\frac{\sin x}{\cos x(\cos x+1)}\,dx.
\]
Let $u=\cos x$, $du=-\sin x\,dx$:
\[
-\int \frac{du}{u(u+1)}=-\int\left(\frac1u-\frac1{u+1}\right)du=\ln\left|\frac{u+1}{u}\right|+C=\ln\left(\frac{\cos x+1}{\cos x}\right)+C.
\]
Evaluating from $0$ to $\pi/3$ (where $\cos x=\tfrac12,\,1$):
\[
\ln\!\left(\frac{3/2}{1/2}\right)-\ln\!\left(\frac{2}{1}\right)=\ln 3-\ln2=\ln\!\left(\frac32\right).
\]
\end{solution}

\question \[\int_{1/2}^1 \frac{dx}{x^2\sqrt{1-x^2}}\]
\begin{solution}
One checks directly that $\dfrac{d}{dx}\left[-\dfrac{\sqrt{1-x^2}}{x}\right]=\dfrac{1}{x^2\sqrt{1-x^2}}$. Hence
\[
\int_{1/2}^1\frac{dx}{x^2\sqrt{1-x^2}}=\left[-\frac{\sqrt{1-x^2}}{x}\right]_{1/2}^1=0-\left(-\frac{\sqrt{3}/2}{1/2}\right)=\sqrt3.
\]
\end{solution}

\question \[\int_2^4 \frac{x^2}{x-1}\,dx\]
\begin{solution}
Polynomial division: $\dfrac{x^2}{x-1}=x+1+\dfrac{1}{x-1}$. So
\[
\int_2^4\left(x+1+\frac{1}{x-1}\right)dx=\left[\frac{x^2}{2}+x+\ln|x-1|\right]_2^4=(12+\ln3)-(4+\ln1)=8+\ln3.
\]
\end{solution}

\end{questions}

\subsubsection*{Regular Round: Bracket 3}
\begin{questions}

\question \[\int_{-2}^2 (x-4)\sqrt{4-x^2}\,dx\]
\begin{solution}
Split into $\displaystyle\int_{-2}^2 x\sqrt{4-x^2}\,dx-4\int_{-2}^2\sqrt{4-x^2}\,dx$. The first integrand is odd on a symmetric interval, so it vanishes. The second integral is the area of a semicircle of radius $2$: $\dfrac12\pi(2)^2=2\pi$. Hence the total is $0-4(2\pi)=-8\pi$.
\end{solution}

\question \[\int_0^\infty \frac{dx}{e^{2x}+e^{-2x}}\]
\begin{solution}
Write $e^{2x}+e^{-2x}=2\cosh(2x)$, so the integral is $\dfrac12\displaystyle\int_0^\infty \operatorname{sech}(2x)\,dx$. With $u=2x$,
\[
\frac12\int_0^\infty \operatorname{sech}(2x)\,dx=\frac14\int_0^\infty \operatorname{sech}(u)\,du=\frac14\Big[\arctan(\sinh u)\Big]_0^\infty=\frac14\cdot\frac{\pi}{2}=\frac{\pi}{8}.
\]
\end{solution}

\question \[\int_0^5 \frac{x+2}{\sqrt{(x+4)x}}\,dx\]
\begin{solution}
Note $\dfrac{d}{dx}(x^2+4x)=2x+4=2(x+2)$, so $(x+2)\,dx=\tfrac12\,d(x^2+4x)$. With $u=x^2+4x$,
\[
\int \frac{x+2}{\sqrt{x^2+4x}}\,dx=\frac12\int \frac{du}{\sqrt u}=\sqrt u=\sqrt{x^2+4x}.
\]
Evaluating from $0$ to $5$: $\sqrt{45}-\sqrt0=3\sqrt5$.
\end{solution}

\end{questions}

\subsubsection*{Regular Round: Bracket 4}
\begin{questions}

\question \[\int_3^4 x\sqrt{x-3}\,dx\]
\begin{solution}
Let $t=x-3$, so $x=t+3$, $dx=dt$, and the limits become $t=0$ to $t=1$:
\[
\int_0^1 (t+3)\sqrt t\,dt=\int_0^1\left(t^{3/2}+3t^{1/2}\right)dt=\left[\frac25 t^{5/2}+2t^{3/2}\right]_0^1=\frac25+2=\frac{12}{5}.
\]
\end{solution}

\question \[\int_0^\pi x^2\cos(x)\,dx\]
\begin{solution}
Integrating by parts twice gives the standard antiderivative
\[
\int x^2\cos x\,dx=x^2\sin x+2x\cos x-2\sin x+C.
\]
Evaluating from $0$ to $\pi$:
\[
(0+2\pi(-1)-0)-(0+0-0)=-2\pi.
\]
\end{solution}

\question \[\int_0^{\pi/3} \tan(x)\ln(\cos(x))\,dx\]
\begin{solution}
Let $u=\ln(\cos x)$, so $du=-\tan x\,dx$, i.e.\ $\tan x\,dx=-du$:
\[
\int \tan x\ln(\cos x)\,dx=\int u(-du)=-\frac{u^2}{2}+C=-\frac12\ln^2(\cos x)+C.
\]
Evaluating from $0$ to $\pi/3$ ($\cos(\pi/3)=\tfrac12$, $\cos0=1$):
\[
-\frac12\ln^2\!\left(\frac12\right)-0=-\frac12\ln^2(2).
\]
\end{solution}

\end{questions}

\subsubsection*{Regular Round: Bracket 5}
\begin{questions}

\question \[\int_1^2 \left(\frac{2}{x+1}+\frac{1}{x(x+1)}\right)dx\]
\begin{solution}
Since $\dfrac{1}{x(x+1)}=\dfrac1x-\dfrac1{x+1}$, the integrand simplifies:
\[
\frac{2}{x+1}+\frac1x-\frac1{x+1}=\frac1x+\frac1{x+1}.
\]
So
\[
\int_1^2\left(\frac1x+\frac1{x+1}\right)dx=\Big[\ln x+\ln(x+1)\Big]_1^2=\ln(2\cdot3)-\ln(1\cdot2)=\ln6-\ln2=\ln3.
\]
\end{solution}

\question \[\int_0^{\pi/4} \frac{1-\tan(x)}{1+\tan(x)}\,dx\]
\begin{solution}
Recall the tangent subtraction formula: $\tan\!\left(\dfrac\pi4-x\right)=\dfrac{1-\tan x}{1+\tan x}$. So
\[
\int_0^{\pi/4}\tan\!\left(\frac\pi4-x\right)dx.
\]
Let $w=\dfrac\pi4-x$: as $x:0\to\pi/4$, $w:\pi/4\to0$, and $dx=-dw$, giving
\[
\int_0^{\pi/4}\tan w\,dw=\Big[-\ln|\cos w|\Big]_0^{\pi/4}=-\ln\!\left(\frac{1}{\sqrt2}\right)+\ln1=\frac12\ln2.
\]
\end{solution}

\question \[\int_0^{\pi/4} \sec^4(x)\,dx\]
\begin{solution}
Write $\sec^4x=\sec^2x(1+\tan^2x)$, and let $u=\tan x$, $du=\sec^2x\,dx$:
\[
\int(1+u^2)\,du=u+\frac{u^3}{3}+C=\tan x+\frac{\tan^3x}{3}+C.
\]
Evaluating from $0$ to $\pi/4$: $\left(1+\dfrac13\right)-0=\dfrac43$.
\end{solution}

\end{questions}

\subsubsection*{Regular Round: Bracket 6}
\begin{questions}

\question \[\int_0^{\sqrt3} \left(\frac{x}{\sqrt{1+x^2}}\right)^3 dx\]
\begin{solution}
Let $u=1+x^2$, $du=2x\,dx$. Write $x^3\,dx=x^2\cdot x\,dx=(u-1)\cdot\dfrac{du}{2}$:
\[
\int \frac{x^3}{(1+x^2)^{3/2}}\,dx=\frac12\int\frac{u-1}{u^{3/2}}\,du=\frac12\int\left(u^{-1/2}-u^{-3/2}\right)du=\sqrt u+\frac1{\sqrt u}+C.
\]
As $x:0\to\sqrt3$, $u:1\to4$. Evaluating: $\left(2+\dfrac12\right)-\left(1+1\right)=\dfrac12$.
\end{solution}

\question \[\int_0^1 \frac{\ln(\sqrt{x})}{\sqrt{x}}\,dx\]
\begin{solution}
Write $\ln\sqrt x=\tfrac12\ln x$, so the integral is $\tfrac12\displaystyle\int_0^1 x^{-1/2}\ln x\,dx$.
Integrate by parts with $u=\ln x,\,dv=x^{-1/2}dx$, so $v=2\sqrt x$:
\[
\int_0^1 x^{-1/2}\ln x\,dx=\Big[2\sqrt x\ln x\Big]_0^1-\int_0^1 2x^{-1/2}\,dx=0-2\Big[2\sqrt x\Big]_0^1=-4.
\]
Multiplying by $\tfrac12$ gives $-2$.
\end{solution}

\question \[\int_0^1 \frac{1-x^{5/2}}{\sqrt{x}}\,dx\]
\begin{solution}
\[
\int_0^1\left(x^{-1/2}-x^2\right)dx=\left[2\sqrt x-\frac{x^3}{3}\right]_0^1=2-\frac13=\frac53.
\]
\end{solution}

\end{questions}

\subsubsection*{Regular Round: Bracket 7}
\begin{questions}

\question \[\int_1^\infty \frac{dx}{\sqrt{e^x-e}}\]
\begin{solution}
Let $u=\sqrt{e^{x-1}-1}$, so $e^{x-1}=1+u^2$, $x=1+\ln(1+u^2)$, and $dx=\dfrac{2u}{1+u^2}\,du$.
Also $e^x-e=e\left(e^{x-1}-1\right)=eu^2$, so $\sqrt{e^x-e}=u\sqrt e$. As $x:1\to\infty$, $u:0\to\infty$. Then
\[
\int_1^\infty \frac{dx}{\sqrt{e^x-e}}=\int_0^\infty \frac{1}{u\sqrt e}\cdot\frac{2u}{1+u^2}\,du=\frac{2}{\sqrt e}\int_0^\infty\frac{du}{1+u^2}=\frac{2}{\sqrt e}\cdot\frac{\pi}{2}=\frac{\pi}{\sqrt e}.
\]
\end{solution}

\question \[\int_0^{\pi/2} \frac{\sin(2x)}{1+\cos^2(x)}\,dx\]
\begin{solution}
Write $\sin2x=2\sin x\cos x$ and let $u=\cos x$, $du=-\sin x\,dx$:
\[
\int \frac{2\sin x\cos x}{1+\cos^2x}\,dx=-2\int\frac{u}{1+u^2}\,du=-\ln(1+u^2)+C.
\]
As $x:0\to\pi/2$, $u:1\to0$. Evaluating $-\ln(1+u^2)$ at $u=0$ minus at $u=1$: $0-(-\ln2)=\ln2$.
\end{solution}

\question \[\int_1^2 \left(x-\frac{1}{x}\right)\left(1+\frac{1}{x^2}\right)dx\]
\begin{solution}
Expand the product:
\[
\left(x-\frac1x\right)\left(1+\frac1{x^2}\right)=x+\frac1x-\frac1x-\frac1{x^3}=x-\frac1{x^3}.
\]
So
\[
\int_1^2\left(x-\frac{1}{x^3}\right)dx=\left[\frac{x^2}{2}+\frac{1}{2x^2}\right]_1^2=\left(2+\frac18\right)-\left(\frac12+\frac12\right)=\frac{17}{8}-1=\frac98.
\]
\end{solution}

\end{questions}

\subsubsection*{Regular Round: Bracket 8}
\begin{questions}

\question \[\int_{3/2}^{5/2} \frac{dx}{x^2-4x+3}\]
\begin{solution}
Factor: $x^2-4x+3=(x-1)(x-3)$, and use partial fractions:
\[
\frac{1}{(x-1)(x-3)}=\frac12\left(\frac{1}{x-3}-\frac{1}{x-1}\right).
\]
An antiderivative is $\dfrac12\ln\left|\dfrac{x-3}{x-1}\right|$. Evaluating from $3/2$ to $5/2$:
at $5/2$: $\left|\dfrac{-1/2}{3/2}\right|=\dfrac13$; at $3/2$: $\left|\dfrac{-3/2}{1/2}\right|=3$.
\[
\frac12\ln\!\left(\frac13\right)-\frac12\ln(3)=\frac12(-\ln3-\ln3)=-\ln3.
\]
\end{solution}

\question \[\int_{1/2}^2 \frac{x}{\sqrt{2x-1}}\,dx\]
\begin{solution}
Let $t=\sqrt{2x-1}$, so $x=\dfrac{t^2+1}{2}$, $dx=t\,dt$. As $x:\tfrac12\to2$, $t:0\to\sqrt3$:
\[
\int_0^{\sqrt3} \frac{(t^2+1)/2}{t}\cdot t\,dt=\frac12\int_0^{\sqrt3}(t^2+1)\,dt=\frac12\left[\frac{t^3}{3}+t\right]_0^{\sqrt3}=\frac12\left(\sqrt3+\sqrt3\right)=\sqrt3.
\]
\end{solution}

\question \[\int_0^\infty \frac{x^2}{1+x^6}\,dx\]
\begin{solution}
Use the standard formula $\displaystyle\int_0^\infty \frac{x^{m-1}}{1+x^n}\,dx=\frac{\pi}{n\sin(m\pi/n)}$ with $m=3$, $n=6$:
\[
\frac{\pi}{6\sin(3\pi/6)}=\frac{\pi}{6\sin(\pi/2)}=\frac{\pi}{6}.
\]
\end{solution}

\end{questions}

\subsubsection*{Quarterfinals: Bracket 1}
\begin{questions}

\question \[\int_0^{\pi/4} \left(\frac{\sin(3x)}{\sin(x)}+\frac{\cos(3x)}{\cos(x)}\right)dx\]
\begin{solution}
Using the triple-angle identities $\sin3x=3\sin x-4\sin^3x$ and $\cos3x=4\cos^3x-3\cos x$:
\[
\frac{\sin3x}{\sin x}=3-4\sin^2x,\qquad \frac{\cos3x}{\cos x}=4\cos^2x-3.
\]
Adding: $3-4\sin^2x+4\cos^2x-3=4(\cos^2x-\sin^2x)=4\cos(2x)$. Hence
\[
\int_0^{\pi/4}4\cos(2x)\,dx=\Big[2\sin(2x)\Big]_0^{\pi/4}=2\sin\!\left(\frac\pi2\right)=2.
\]
\end{solution}

\question \[\int e^{x+\frac1x}\left(x+1-\frac1x\right)dx\]
\begin{solution}
Guess the antiderivative $F(x)=xe^{x+1/x}$ and differentiate:
\[
F'(x)=e^{x+1/x}+x e^{x+1/x}\left(1-\frac{1}{x^2}\right)=e^{x+1/x}\left(1+x-\frac1x\right),
\]
which is exactly the integrand. Hence
\[
\int e^{x+\frac1x}\left(x+1-\frac1x\right)dx=xe^{x+\frac1x}+C.
\]
\end{solution}

\question \[\int_0^1 \left(\sqrt{1+9x^4}-\sqrt{1+\left(\frac{1}{9x^4}\right)^{1/3}}+x^4\right)dx\]
\begin{solution}
The first square root, $\sqrt{1+9x^4}=\sqrt{1+(3x^2)^2}$, is the arclength density (speed) of the curve $y=x^3$, since $y'=3x^2$. The second square-root term is exactly the arclength density of the \emph{inverse} curve $x=y^{1/3}$ traced over the same interval. Because a curve and its reflection across $y=x$ (its inverse) have equal arclength when both are traversed between the same two endpoints — here $(0,0)$ and $(1,1)$ — the two square-root integrals are equal:
\[
\int_0^1\sqrt{1+9x^4}\,dx=\int_0^1\sqrt{1+\left(\frac1{9x^4}\right)^{1/3}}dx.
\]
So they cancel, leaving only
\[
\int_0^1 x^4\,dx=\frac15.
\]
\end{solution}

\end{questions}

\subsubsection*{Quarterfinals: Bracket 2}
\begin{questions}

\question \[\int_{\pi/6}^{\pi/3} \frac{dx}{(\cos(x)-1)^2}\]
\begin{solution}
Since $1-\cos x=2\sin^2(x/2)$, we get $(\cos x-1)^2=4\sin^4(x/2)$, so the integrand is $\tfrac14\csc^4(x/2)$.
Using $\displaystyle\int\csc^4u\,du=-\cot u-\frac{\cot^3u}{3}+C$ with $u=x/2$ (so $du=dx/2$):
\[
\int \csc^4(x/2)\,dx=2\left(-\cot\frac x2-\frac{\cot^3(x/2)}{3}\right)+C.
\]
Thus an antiderivative is $F(x)=-\dfrac12\cot\dfrac x2-\dfrac16\cot^3\dfrac x2$.
At $x=\pi/3$: $\cot(\pi/6)=\sqrt3$, giving $F=-\dfrac{\sqrt3}{2}-\dfrac{3\sqrt3}{6}=-\sqrt3$.
At $x=\pi/6$: $\cot(\pi/12)=2+\sqrt3$, and $(2+\sqrt3)^3=26+15\sqrt3$, giving
\[
F=-\frac{2+\sqrt3}{2}-\frac{26+15\sqrt3}{18}=-\frac{16}{3}-3\sqrt3\quad\text{(after simplifying)}.
\]
Subtracting: $F(\pi/3)-F(\pi/6)=-\sqrt3-\left(-\dfrac{16}{3}-3\sqrt3\right)=\dfrac{16}{3}+2\sqrt3$.
\end{solution}

\question \[\int_{-\infty}^{\infty} \frac{12^x}{9^x+16^x}\,dx\]
\begin{solution}
Divide numerator and denominator by $16^x$: since $12/16=3/4$ and $9/16=(3/4)^2$,
\[
\int \frac{(3/4)^x}{1+(3/4)^{2x}}\,dx.
\]
Let $t=(3/4)^x$, so $dt=\ln(3/4)\,t\,dx$. Then
\[
\int \frac{t}{1+t^2}\cdot\frac{dt}{t\ln(3/4)}=\frac{1}{\ln(3/4)}\arctan(t)+C.
\]
As $x:-\infty\to\infty$, $t=(3/4)^x$ goes from $\infty$ to $0$ (since $3/4<1$). Hence
\[
\frac{1}{\ln(3/4)}\Big[\arctan(0)-\arctan(\infty)\Big]=\frac{-\pi/2}{\ln(3/4)}=\frac{\pi/2}{\ln(4/3)}=\frac{\pi}{2\ln(4/3)}.
\]
\end{solution}

\question \[\int_0^\pi \Big(\sin(x)\sin(2x)\sin(4x)+\cos(x)\cos(2x)\cos(4x)\Big)\,dx\]
\begin{solution}
Expand each triple product using product-to-sum formulas repeatedly:
\[
\sin x\sin2x\sin4x=\frac14\big(\sin3x-\sin x-\sin7x+\sin5x\big),
\]
\[
\cos x\cos2x\cos4x=\frac14\big(\cos x+\cos3x+\cos5x+\cos7x\big).
\]
Adding,
\[
\sin x\sin2x\sin4x+\cos x\cos2x\cos4x=\frac14\Big[(\cos x-\sin x)+(\cos3x+\sin3x)+(\cos5x+\sin5x)+(\cos7x-\sin7x)\Big].
\]
Now use $\displaystyle\int_0^\pi\cos(kx)\,dx=0$ and $\displaystyle\int_0^\pi\sin(kx)\,dx=\frac{2}{k}$ for odd integers $k$. So the integral equals
\[
\frac14\left[-\frac21+\frac23+\frac25-\frac27\right]=\frac12\left[-1+\frac13+\frac15-\frac17\right]=\frac12\cdot\left(-\frac{64}{105}\right)=-\frac{32}{105}.
\]
\end{solution}

\end{questions}

\subsubsection*{Quarterfinals: Bracket 3}
\begin{questions}

\question \[\int_0^1 \sqrt{\frac{1}{1-x}+\frac{x-1}{x^2+x+1}}\,dx\]
\begin{solution}
Combine the two fractions over the common denominator $(1-x)(x^2+x+1)=1-x^3$:
\[
\frac{1}{1-x}-\frac{1-x}{x^2+x+1}=\frac{(x^2+x+1)-(1-x)^2}{1-x^3}=\frac{3x}{1-x^3}.
\]
So the integrand is $\sqrt{\dfrac{3x}{1-x^3}}$. Substituting $x^3=t$, i.e. $x=t^{1/3}$, $dx=\tfrac13 t^{-2/3}dt$:
\[
\int_0^1\sqrt3\,x^{1/2}(1-x^3)^{-1/2}dx=\frac{\sqrt3}{3}\int_0^1 t^{-1/2}(1-t)^{-1/2}\,dt=\frac{\sqrt3}{3}B\!\left(\frac12,\frac12\right)=\frac{\sqrt3}{3}\pi=\frac{\pi}{\sqrt3}.
\]
\end{solution}

\question \[\int_0^\infty \left(\frac{1}{\sqrt{x^2+\pi^2}}-\frac{1}{x+\pi}\right)dx\]
\begin{solution}
For any constant $a>0$, an antiderivative of the difference is
\[
F(x)=\ln\!\left(x+\sqrt{x^2+a^2}\right)-\ln(x+a).
\]
At $x=0$: $F(0)=\ln(a)-\ln(a)=0$. As $x\to\infty$, $\sqrt{x^2+a^2}\sim x$, so $x+\sqrt{x^2+a^2}\sim 2x$ while $x+a\sim x$, giving $F(\infty)=\ln2$.
So the (finite) value is $\ln2-0=\ln2$, independent of $a$; here $a=\pi$.
\end{solution}

\question \[\int_0^1 \frac{\sqrt{1-x^2}}{(1+x)^2}\,dx\]
\begin{solution}
Write $\sqrt{1-x^2}=\sqrt{1-x}\sqrt{1+x}$, so the integrand is $\dfrac{\sqrt{1-x}}{(1+x)^{3/2}}$.
Let $t=\sqrt{\dfrac{1-x}{1+x}}$, so $x=\dfrac{1-t^2}{1+t^2}$, $1+x=\dfrac{2}{1+t^2}$, and $dx=-\dfrac{4t}{(1+t^2)^2}\,dt$. Then
\[
\frac{\sqrt{1-x}}{(1+x)^{3/2}}\,dx = t\cdot\frac{1+t^2}{2}\cdot\left(-\frac{4t}{(1+t^2)^2}\right)dt=-\frac{2t^2}{1+t^2}\,dt.
\]
As $x:0\to1$, $t:1\to0$, so
\[
\int_0^1 \frac{2t^2}{1+t^2}\,dt=2\int_0^1\left(1-\frac{1}{1+t^2}\right)dt=2\Big[t-\arctan t\Big]_0^1=2\left(1-\frac{\pi}{4}\right)=2-\frac{\pi}{2}.
\]
\end{solution}

\end{questions}

\subsubsection*{Quarterfinals: Bracket 4}
\begin{questions}

\question \[\int_{\pi/4}^{\pi/3} \frac{(\sec(x)-\csc(x))(\sec(x)+\csc(x))}{\tan(x)+\cot(x)}\,dx\]
\begin{solution}
The numerator is $\sec^2x-\csc^2x$. The denominator simplifies:
\[
\tan x+\cot x=\frac{\sin x}{\cos x}+\frac{\cos x}{\sin x}=\frac{1}{\sin x\cos x}.
\]
So the integrand is $(\sec^2x-\csc^2x)\sin x\cos x=\dfrac{\sin x}{\cos x}-\dfrac{\cos x}{\sin x}=\tan x-\cot x$.
An antiderivative is $-\ln|\cos x|-\ln|\sin x|=-\ln|\sin x\cos x|$.
Since $\sin x\cos x=\tfrac12\sin2x$: at $x=\pi/3$, $\sin x\cos x=\tfrac{\sqrt3}{4}$; at $x=\pi/4$, $\sin x \cos x=\tfrac12$.
\[
-\ln\!\left(\frac{\sqrt3}{4}\right)+\ln\!\left(\frac12\right)=\ln\!\left(\frac{4}{\sqrt3}\cdot\frac12\right)=\ln\!\left(\frac{2}{\sqrt3}\right)=\ln2-\frac12\ln3=\frac12\ln\!\left(\frac43\right).
\]
\end{solution}

\question \[\int_1^\infty \frac{dx}{x\sqrt{x}\sqrt{x-1}}\]
\begin{solution}
Let $x-1=u^2$, so $x=1+u^2$, $dx=2u\,du$, $\sqrt{x-1}=u$:
\[
\int_0^\infty \frac{2u\,du}{(1+u^2)^{3/2}\,u}=2\int_0^\infty \frac{du}{(1+u^2)^{3/2}}=2\left[\frac{u}{\sqrt{1+u^2}}\right]_0^\infty=2(1-0)=2.
\]
\end{solution}

\question \[\int \frac{\tan(x)+\cot(x)}{\ln(\tan(x))}\,dx\]
\begin{solution}
Let $u=\ln(\tan x)$. Then
\[
du=\frac{\sec^2x}{\tan x}\,dx=\frac{1}{\sin x\cos x}\,dx=(\tan x+\cot x)\,dx.
\]
So the integral becomes $\displaystyle\int \frac{du}{u}=\ln|u|+C=\ln\big|\ln(\tan x)\big|+C$.
\end{solution}

\end{questions}

\subsubsection*{Semifinals: Bracket 1}
\begin{questions}

\question \[\int_0^\infty \frac{\arctan(x)}{(x+1)^3}\,dx\]
\begin{solution}
Integrate by parts with $u=\arctan x$, $dv=(x+1)^{-3}dx$, $v=-\dfrac{1}{2(x+1)^2}$:
\[
\int_0^\infty \frac{\arctan x}{(x+1)^3}dx=\left[-\frac{\arctan x}{2(x+1)^2}\right]_0^\infty+\frac12\int_0^\infty \frac{dx}{(x+1)^2(1+x^2)}.
\]
The boundary term vanishes at both ends. Using partial fractions,
\[
\frac{1}{(x+1)^2(1+x^2)}=\frac12\left[\frac{1}{x+1}+\frac{1}{(x+1)^2}-\frac{x}{1+x^2}\right].
\]
Now $\displaystyle\int_0^\infty\left[\frac1{x+1}-\frac{x}{1+x^2}\right]dx=\left[\ln\frac{x+1}{\sqrt{1+x^2}}\right]_0^\infty=0-0=0$ (the ratio tends to $1$ at infinity), while $\displaystyle\int_0^\infty\frac{dx}{(x+1)^2}=1$.
So the remaining piece is $\dfrac12\cdot\dfrac12(0+1)=\dfrac14$.
\end{solution}

\question \[\int_0^1 \frac{9x^2}{x\sqrt{x}+\sqrt{x^3+1}}\,dx\]
\begin{solution}
Rationalize by multiplying by $\dfrac{\sqrt{x^3+1}-x^{3/2}}{\sqrt{x^3+1}-x^{3/2}}$; the denominator becomes $(x^3+1)-x^3=1$. So
\[
\frac{1}{x^{3/2}+\sqrt{x^3+1}}=\sqrt{x^3+1}-x^{3/2}.
\]
The integrand becomes $9x^2\sqrt{x^3+1}-9x^{7/2}$. For the first term let $w=x^3+1$, $dw=3x^2dx$:
\[
\int 9x^2\sqrt{x^3+1}\,dx=3\int\sqrt w\,dw=2w^{3/2}=2(x^3+1)^{3/2}.
\]
For the second: $\displaystyle\int 9x^{7/2}dx=2x^{9/2}$. So an antiderivative is $2(x^3+1)^{3/2}-2x^{9/2}$.
Evaluating from $0$ to $1$: $\big(2\cdot2^{3/2}-2\big)-\big(2-0\big)=4\sqrt2-2-2=4\sqrt2-4$.
\end{solution}

\question \[\int_{-\pi}^\pi \frac{dx}{(1+e^{\cos(x)})(1+e^{x\cos(x)})}\]
\begin{solution}
Let $f(x)$ denote the integrand, $I=\int_{-\pi}^\pi f(x)\,dx$. Note $\cos x$ is even and $x\cos x$ is odd. Using the identity $\dfrac{1}{1+e^t}+\dfrac{1}{1+e^{-t}}=1$ with $t=x\cos x$:
\[
f(x)+f(-x)=\frac{1}{1+e^{\cos x}}\left(\frac{1}{1+e^{x\cos x}}+\frac{1}{1+e^{-x\cos x}}\right)=\frac{1}{1+e^{\cos x}}.
\]
So $2I=\displaystyle\int_{-\pi}^\pi \frac{dx}{1+e^{\cos x}}=:J$.
Since the integrand of $J$ has period $2\pi$, we may also write $J=\int_0^{2\pi}\dfrac{dy}{1+e^{\cos y}}$. Pairing $y$ with $y+\pi$ (using $\cos(y+\pi)=-\cos y$ and the same identity above):
\[
\frac{1}{1+e^{\cos y}}+\frac{1}{1+e^{\cos(y+\pi)}}=\frac{1}{1+e^{\cos y}}+\frac{1}{1+e^{-\cos y}}=1.
\]
So $J=\displaystyle\int_0^\pi 1\,dy=\pi$. Hence $2I=\pi$, so $I=\pi/2$.
\end{solution}

\end{questions}

\subsubsection*{Semifinals: Bracket 2}
\begin{questions}

\question \[\int_2^5 \frac{dx}{\sqrt{x-2\sqrt{x-2\sqrt{x-2\sqrt{\cdots}}}}}\]
\begin{solution}
Let $L=L(x)=\sqrt{x-2\sqrt{x-2\sqrt{x-\cdots}}}$. By self-similarity, $L=\sqrt{x-2L}$, so $L^2+2L-x=0$, giving (taking the nonnegative root) $L=\sqrt{x+1}-1$.
So the integrand is $\dfrac{1}{\sqrt{x+1}-1}$. Let $t=\sqrt{x+1}$, $x=t^2-1$, $dx=2t\,dt$:
\[
\int \frac{2t}{t-1}\,dt=2\int\left(1+\frac{1}{t-1}\right)dt=2t+2\ln|t-1|+C=2\sqrt{x+1}+2\ln\big(\sqrt{x+1}-1\big)+C.
\]
Evaluating from $x=2$ ($t=\sqrt3$) to $x=5$ ($t=\sqrt6$):
\[
\Big(2\sqrt6+2\ln(\sqrt6-1)\Big)-\Big(2\sqrt3+2\ln(\sqrt3-1)\Big)=2\sqrt6-2\sqrt3+2\ln\!\left(\frac{\sqrt6-1}{\sqrt3-1}\right).
\]
\end{solution}

\question \[\int_0^{\pi/3} \sqrt{\cos(x)}\left(\sec^2(x)+1\right)dx\]
\begin{solution}
Consider $h(x)=\dfrac{2\sin x}{\sqrt{\cos x}}$. Differentiating,
\[
h'(x)=2\sqrt{\cos x}+2\sin x\cdot\left(-\frac12\right)\cos^{-3/2}x\cdot(-\sin x)=2\sqrt{\cos x}+\sin^2x\cos^{-3/2}x.
\]
Since $\sin^2x=1-\cos^2x$, the second term is $\cos^{-3/2}x-\sqrt{\cos x}$, so
\[
h'(x)=\sqrt{\cos x}+\cos^{-3/2}x=\sqrt{\cos x}\left(\sec^2x+1\right),
\]
exactly the integrand. So
\[
\int_0^{\pi/3}\sqrt{\cos x}(\sec^2x+1)\,dx=\Big[h(x)\Big]_0^{\pi/3}=\frac{2\cdot\frac{\sqrt3}{2}}{\sqrt{1/2}}-0=\sqrt3\cdot\sqrt2=\sqrt6.
\]
\end{solution}

\question \[\int_{\pi/3}^{\pi/2} \frac{dx}{\sin(x)+\tan(x)}\]
\begin{solution}
Write $\sin x+\tan x=\dfrac{\sin x(\cos x+1)}{\cos x}$, so the integrand is $\dfrac{\cos x}{\sin x(1+\cos x)}$.
Substitute $t=\cos x$, $dt=-\sin x\,dx$, so $dx=-\dfrac{dt}{\sqrt{1-t^2}}$, and $\sin x=\sqrt{1-t^2}$:
\[
\frac{\cos x}{\sin x(1+\cos x)}\,dx=-\frac{t\,dt}{(1-t^2)(1+t)}=-\frac{t\,dt}{(1-t)(1+t)^2}.
\]
Partial fractions give
\[
-\frac{t}{(1-t)(1+t)^2}=-\frac{1/4}{1-t}-\frac{1/4}{1+t}+\frac{1/2}{(1+t)^2}.
\]
Integrating, $\displaystyle\int = \frac14\ln\!\left(\frac{1-t}{1+t}\right)-\frac{1}{2(1+t)}+C$, i.e.\ with $t=\cos x$:
\[
G(x)=\frac14\ln\!\left(\frac{1-\cos x}{1+\cos x}\right)-\frac{1}{2(1+\cos x)}.
\]
At $x=\pi/2$ ($\cos x=0$): $G=0-\dfrac12=-\dfrac12$.
At $x=\pi/3$ ($\cos x=\tfrac12$): $G=\dfrac14\ln\!\left(\dfrac13\right)-\dfrac13=-\dfrac14\ln3-\dfrac13$.
\[
G(\pi/2)-G(\pi/3)=-\frac12-\left(-\frac14\ln3-\frac13\right)=\frac14\ln3-\frac16.
\]
\end{solution}

\end{questions}

\subsubsection*{3rd Place Match}
\begin{questions}

\question \[\int_0^{\pi/2} \frac{dx}{(x\tan(x)+1)^2}\]
\begin{solution}
Write $1+x\tan x=\dfrac{\cos x+x\sin x}{\cos x}$, so the integrand is $\dfrac{\cos^2 x}{(\cos x+x\sin x)^2}$.
Let $D(x)=\cos x+x\sin x$; then $D'(x)=-\sin x+\sin x+x\cos x=x\cos x$.
Consider $\dfrac{d}{dx}\left[\dfrac{\sin x}{D(x)}\right]=\dfrac{\cos x\, D-\sin x\, D'}{D^2}=\dfrac{\cos x(\cos x+x\sin x)-\sin x(x\cos x)}{D^2}=\dfrac{\cos^2x}{D^2}$,
exactly the integrand. So
\[
\int_0^{\pi/2}\frac{dx}{(1+x\tan x)^2}=\left[\frac{\sin x}{\cos x+x\sin x}\right]_0^{\pi/2}=\frac{1}{\pi/2}-0=\frac{2}{\pi}.
\]
\end{solution}

\question \[\int_0^1 \Big(\ln(1-\sqrt[7]{x})-\ln(1-\sqrt[8]{x})\Big)\,dx\]
\begin{solution}
For a positive integer $n$, substitute $x=u^n$ in $\displaystyle\int_0^1\ln(1-x^{1/n})\,dx$:
\[
\int_0^1\ln(1-u)\,nu^{n-1}\,du=n\int_0^1 u^{n-1}\ln(1-u)\,du.
\]
Using the classical result $\displaystyle\int_0^1 u^{n-1}\ln(1-u)\,du=-\frac{H_n}{n}$ (where $H_n=1+\tfrac12+\cdots+\tfrac1n$), this equals $-H_n$. So
\[
\int_0^1\ln\big(1-x^{1/n}\big)\,dx=-H_n.
\]
Therefore
\[
\int_0^1\Big(\ln(1-x^{1/7})-\ln(1-x^{1/8})\Big)dx=-H_7-(-H_8)=H_8-H_7=\frac18.
\]
\end{solution}

\question \[\int_1^\infty \frac{dx}{x\sqrt{x^2-1}\left(1+\sqrt{x^2-1}\right)\left(x^2-\sqrt{x^2-1}\right)}\]
\begin{solution}
Substitute $x=\sec\theta$, so $\sqrt{x^2-1}=\tan\theta$, $dx=\sec\theta\tan\theta\,d\theta$, $\theta:0\to\pi/2$. The denominator becomes
\[
\sec\theta\tan\theta(1+\tan\theta)(\sec^2\theta-\tan\theta),
\]
so after cancelling $\sec\theta\tan\theta$ with $dx$'s factor,
\[
\frac{dx}{(\cdots)}=\frac{d\theta}{(1+\tan\theta)(\sec^2\theta-\tan\theta)}=\frac{d\theta}{(1+\tan\theta)(1+\tan^2\theta-\tan\theta)}=\frac{d\theta}{1+\tan^3\theta},
\]
using $(1+t)(1-t+t^2)=1+t^3$ with $t=\tan\theta$.
So $I=\displaystyle\int_0^{\pi/2}\frac{d\theta}{1+\tan^3\theta}$. Substituting $\theta\to\pi/2-\theta$ sends $\tan\theta\to\cot\theta$, giving
\[
I=\int_0^{\pi/2}\frac{d\theta}{1+\cot^3\theta}=\int_0^{\pi/2}\frac{\tan^3\theta}{1+\tan^3\theta}\,d\theta.
\]
Adding this to the original expression for $I$:
\[
2I=\int_0^{\pi/2}\left[\frac{1}{1+\tan^3\theta}+\frac{\tan^3\theta}{1+\tan^3\theta}\right]d\theta=\int_0^{\pi/2}1\,d\theta=\frac\pi2,
\]
so $I=\pi/4$.
\end{solution}

\question \[\int_{-\infty}^\infty \frac{e^x}{\sinh^2(x)+2}\,dx\]
\begin{solution}
Since $\cosh^2x-\sinh^2x=1$, we have $\sinh^2x+2=\cosh^2x+1$. Writing $e^x=\cosh x+\sinh x$, split the integral:
\[
\int_{-\infty}^\infty \frac{\cosh x}{\cosh^2x+1}\,dx+\int_{-\infty}^\infty \frac{\sinh x}{\cosh^2x+1}\,dx.
\]
For the second integral, let $u=\cosh x$: it becomes $\displaystyle\int \frac{du}{u^2+1}=\arctan(\cosh x)$; since $\cosh x\to\infty$ as $x\to\pm\infty$ symmetrically, this integral evaluates to $0$.
For the first, use $\cosh^2x+1=2+\sinh^2x$ and let $v=\sinh x$, $dv=\cosh x\,dx$:
\[
\int \frac{dv}{v^2+2}=\frac{1}{\sqrt2}\arctan\!\left(\frac{v}{\sqrt2}\right)=\frac{1}{\sqrt2}\arctan\!\left(\frac{\sinh x}{\sqrt2}\right).
\]
As $x\to\pm\infty$, $\sinh x\to\pm\infty$, so this contributes $\dfrac{1}{\sqrt2}\left(\dfrac\pi2-\left(-\dfrac\pi2\right)\right)=\dfrac{\pi}{\sqrt2}$.
Total: $\dfrac{\pi}{\sqrt2}+0=\dfrac{\pi}{\sqrt2}$.
\end{solution}

\question \[\int_0^\infty \frac{dx}{\sqrt{x\sqrt{x}}\left(x^{1/4}+1\right)^{10}}\]
\begin{solution}
Note $\sqrt{x\sqrt x}=\sqrt{x^{3/2}}=x^{3/4}$. Let $t=x^{1/4}$, so $x=t^4$, $dx=4t^3\,dt$, and $x^{3/4}=t^3$:
\[
\frac{dx}{x^{3/4}(t+1)^{10}}=\frac{4t^3\,dt}{t^3(t+1)^{10}}=\frac{4\,dt}{(t+1)^{10}}.
\]
As $x:0\to\infty$, $t:0\to\infty$:
\[
\int_0^\infty \frac{4\,dt}{(t+1)^{10}}=4\left[-\frac{1}{9(t+1)^9}\right]_0^\infty=4\cdot\frac19=\frac49.
\]
\end{solution}

\end{questions}

\subsubsection*{Finals Round}
\begin{questions}

\question \[\sum_{n=1}^{2024}\int_{1/2^{n+1}}^{1/2^n} \frac{\ln(2^n x)\ln(2^{n+1}x)\ln(2^{n+2}x)}{x}\,dx\]
\begin{solution}
Fix $n$ and substitute $u=\ln(2^nx)=n\ln2+\ln x$, so $du=\dfrac{dx}{x}$. As $x$ runs from $2^{-(n+1)}$ to $2^{-n}$, $u$ runs from $\ln(2^n\cdot2^{-(n+1)})=-\ln2$ to $\ln(2^n\cdot2^{-n})=0$. Also $\ln(2^{n+1}x)=u+\ln2$ and $\ln(2^{n+2}x)=u+2\ln2$. Writing $a=\ln2$, the $n$-th term becomes
\[
\int_{-a}^{0} u(u+a)(u+2a)\,du,
\]
which is \emph{independent of $n$}. Expanding, $u(u+a)(u+2a)=u^3+3au^2+2a^2u$, with antiderivative $\dfrac{u^4}{4}+au^3+a^2u^2$. Evaluating from $-a$ to $0$:
\[
0-\left(\frac{a^4}{4}-a^4+a^4\right)=0-\frac{a^4}{4}=-\frac{a^4}{4}.
\]
So every one of the $2024$ terms equals $-\dfrac{(\ln2)^4}{4}$, and the sum is
\[
2024\cdot\left(-\frac{(\ln2)^4}{4}\right)=-506(\ln2)^4.
\]
\end{solution}

\question \[\int_0^{\pi/2} \frac{\arctan(2\tan(x))}{\tan(x)}\,dx\]
\begin{solution}
Define $I(a)=\displaystyle\int_0^{\pi/2}\frac{\arctan(a\tan x)}{\tan x}\,dx$, so $I(0)=0$ and we want $I(2)$. Differentiate under the integral sign:
\[
I'(a)=\int_0^{\pi/2}\frac{1}{1+a^2\tan^2x}\,dx.
\]
This is a standard integral with value $\dfrac{\pi}{2(1+a)}$ for $a\ge0$ (obtained, e.g., via the Weierstrass substitution). Hence
\[
I(a)=\int_0^a \frac{\pi}{2(1+s)}\,ds=\frac{\pi}{2}\ln(1+a).
\]
At $a=2$: $I(2)=\dfrac{\pi}{2}\ln3$.
\end{solution}

\question \[\lim_{n\to\infty}\frac1n\int_{-n}^{n}\frac{dx}{e^{-x^x}+1}\]
\begin{solution}
Let $f(x)=\dfrac{1}{1+e^{-x^x}}$, which always satisfies $0<f(x)<1$.
For $x>0$ large, $x^x\to\infty$ extremely quickly, so $e^{-x^x}\to0$ and $f(x)\to1$; only a bounded region near $x=0$ fails to be close to $1$, which contributes a vanishing fraction as $n\to\infty$. Hence
\[
\frac1n\int_0^n f(x)\,dx\longrightarrow 1.
\]
For $x<0$, $x^x$ does not settle to a definite limiting sign or magnitude, so $f(x)$ oscillates boundedly between $0$ and $1$ rather than converging; by an equidistribution/averaging argument its Ces\`aro (running) average tends to $\tfrac12$. Hence
\[
\frac1n\int_{-n}^0 f(x)\,dx\longrightarrow \frac12.
\]
Adding the two pieces,
\[
\lim_{n\to\infty}\frac1n\int_{-n}^n f(x)\,dx=1+\frac12=\frac32.
\]
\end{solution}

\question \[\int_0^{1/2} \ln\left(\prod_{n=1}^{\infty}\sum_{k=0}^{\infty}\frac{n^k}{x^{nk}\cdot k!}\right)dx\]
\begin{solution}
For fixed $n$, the inner sum is an exponential series:
\[
\sum_{k=0}^\infty \frac{n^k}{x^{nk}k!}=\sum_{k=0}^\infty \frac{(n/x^n)^k}{k!}=\exp\!\left(\frac{n}{x^n}\right).
\]
So the product is $\exp\!\left(\sum_{n=1}^\infty n x^{-n}\right)$, and taking the logarithm gives $\sum_{n=1}^\infty n x^{-n}$. Using the generating function $\sum_{n=1}^\infty n r^n=\dfrac{r}{(1-r)^2}$ (applied formally with $r=1/x$),
\[
\sum_{n=1}^\infty n x^{-n}=\frac{1/x}{\left(1-1/x\right)^2}=\frac{x}{(x-1)^2}.
\]
So we must evaluate $\displaystyle\int_0^{1/2}\frac{x}{(x-1)^2}\,dx$. With $u=x-1$: $\dfrac{x}{(x-1)^2}=\dfrac{u+1}{u^2}=\dfrac1u+\dfrac1{u^2}$, so an antiderivative is $\ln|x-1|-\dfrac1{x-1}$.
Evaluating from $0$ to $1/2$:
\[
\left[\ln\!\left(\frac12\right)-\frac{1}{-1/2}\right]-\left[\ln(1)-\frac{1}{-1}\right]=(-\ln2+2)-(0+1)=1-\ln2.
\]
\end{solution}

\question \[\int_0^\infty \frac{\arctan(x)}{(2x+3)(3x+2)}\,dx\]
\begin{solution}
Partial fractions give
\[
\frac{1}{(2x+3)(3x+2)}=-\frac{2/5}{2x+3}+\frac{3/5}{3x+2}.
\]
Introduce a parameter and consider $I(a)=\displaystyle\int_0^\infty \arctan(ax)\left(-\frac{2/5}{2x+3}+\frac{3/5}{3x+2}\right)dx$, so $I(0)=0$ and the desired value is $I(1)$. As $x\to\infty$ the bracket behaves like $O(1/x^2)$, which is enough for $I(a)$ and its derivative to converge for $a>0$. Differentiating under the integral sign and evaluating the resulting rational-function integrals (each reducing, after partial fractions in $x$, to elementary arctangent and logarithm antiderivatives) and integrating back over $a$ from $0$ to $1$ yields the closed form
\[
I(1)=\frac{\pi}{10}\ln\!\left(\frac32\right).
\]
\end{solution}

\end{questions}

\subsubsection*{Finals Tiebreakers}
\begin{questions}

\question \[\int_1^{\int_1^{\ddots}\frac{dx}{\sqrt{x-1}}}\frac{dx}{\sqrt{x-1}}\]
\begin{solution}
Let $L$ denote the value of this self-referential (infinitely nested) integral, so that $L$ is defined by the fixed-point equation
\[
L=\int_1^L \frac{dx}{\sqrt{x-1}}=\Big[2\sqrt{x-1}\Big]_1^L=2\sqrt{L-1}.
\]
Squaring: $L^2=4(L-1)$, i.e.\ $L^2-4L+4=0$, i.e.\ $(L-2)^2=0$. So $L=2$.
\end{solution}

\question \[\int_0^\infty (x+1)^4 e^{-x}\,dx\]
\begin{solution}
Expand $(x+1)^4=x^4+4x^3+6x^2+4x+1$ and use $\displaystyle\int_0^\infty x^n e^{-x}\,dx=n!$:
\[
4!+4\cdot3!+6\cdot2!+4\cdot1!+0!=24+24+12+4+1=65.
\]
\end{solution}

\question \[\int_0^{\pi/2} \sqrt{1-\cos(x)}\sqrt{1-\cos(2x)}\,dx\]
\begin{solution}
Using $1-\cos x=2\sin^2(x/2)$ and $1-\cos2x=2\sin^2x$, and that all quantities are nonnegative on $(0,\pi/2)$:
\[
\sqrt{1-\cos x}\,\sqrt{1-\cos2x}=\sqrt2\sin\!\left(\frac x2\right)\cdot\sqrt2\sin x=2\sin\!\left(\frac x2\right)\sin x.
\]
Since $\sin x=2\sin(x/2)\cos(x/2)$, this is $4\sin^2(x/2)\cos(x/2)$. Let $u=\sin(x/2)$, $du=\tfrac12\cos(x/2)\,dx$:
\[
\int 4\sin^2\!\left(\frac x2\right)\cos\!\left(\frac x2\right)dx=8\int u^2\,du=\frac{8u^3}{3}=\frac83\sin^3\!\left(\frac x2\right).
\]
Evaluating from $0$ to $\pi/2$ ($\sin(\pi/4)=\tfrac{\sqrt2}{2}$):
\[
\frac83\left(\frac{\sqrt2}{2}\right)^3-0=\frac83\cdot\frac{2\sqrt2}{8}=\frac{2\sqrt2}{3}.
\]
\end{solution}

\question \[\int_{-1}^0 \frac{x}{1+\sqrt{x+1}}\,dx\]
\begin{solution}
Let $t=\sqrt{x+1}$, so $x=t^2-1$, $dx=2t\,dt$; as $x:-1\to0$, $t:0\to1$:
\[
\int_0^1 \frac{t^2-1}{1+t}\cdot2t\,dt=\int_0^1 \frac{(t-1)(t+1)}{1+t}\cdot2t\,dt=\int_0^1 2t(t-1)\,dt=\int_0^1(2t^2-2t)\,dt.
\]
\[
=\left[\frac{2t^3}{3}-t^2\right]_0^1=\frac23-1=-\frac13.
\]
\end{solution}

\question \[\int_{-1}^1 (x^2+x+1)^2\,dx\]
\begin{solution}
Expand: $(x^2+x+1)^2=x^4+2x^3+3x^2+2x+1$. On the symmetric interval $[-1,1]$, the odd-degree terms ($2x^3$ and $2x$) integrate to $0$, leaving
\[
\int_{-1}^1 (x^4+3x^2+1)\,dx=2\int_0^1(x^4+3x^2+1)\,dx=2\left(\frac15+1+1\right)=2\cdot\frac{11}{5}=\frac{22}{5}.
\]
\end{solution}

\end{questions}